% --- PHỤ LỤC ---
\appendix

\chapter{Chứng minh toán học bổ trợ}

Phụ lục này trình bày các kết quả toán học nền tảng được sử dụng trong phân tích lý thuyết của nghiên cứu. Mỗi kết quả được giải thích về mặt vai trò và ý nghĩa đối với các chứng minh chính.

\section{Các kết quả đại số tuyến tính}

\subsection{Định lý Sherman-Morrison}

Định lý \textit{Sherman-Morrison} (1950) là công cụ quan trọng để tính nghịch đảo của ma trận khi có sự thay đổi hạng 1 (rank-1 update). Đây là nền tảng cho việc phân tích ma trận NTK có cấu trúc ``scaled identity + rank-1 noise''.

\textbf{Định lý (Sherman-Morrison):} Cho $A$ là ma trận khả nghịch kích thước $n \times n$ và $u, v$ là các vector cột kích thước $n$. Nếu $1 + v^\top A^{-1} u \neq 0$, thì:
\begin{equation}
    (A + uv^\top)^{-1} = A^{-1} - \frac{A^{-1} u v^\top A^{-1}}{1 + v^\top A^{-1} u}
\end{equation}

\textbf{Vai trò trong nghiên cứu:}

Ma trận NTK trên tập huấn luyện có dạng:
\begin{equation}
    \Theta = \alpha I_M + \beta \mathbf{1}\mathbf{1}^\top + \text{(nhiễu bậc cao)}
\end{equation}

trong đó $\mathbf{1}\mathbf{1}^\top$ là ma trận rank-1. Áp dụng Sherman-Morrison với $A = \alpha I_M$, $u = v = \sqrt{\beta}\mathbf{1}$:
\begin{align}
    \Theta^{-1} &\approx (\alpha I_M + \beta \mathbf{1}\mathbf{1}^\top)^{-1} \\
    &= \frac{1}{\alpha} I_M - \frac{\frac{1}{\alpha^2} \beta \mathbf{1}\mathbf{1}^\top}{1 + \frac{\beta M}{\alpha}} \\
    &= \frac{1}{\alpha} I_M - \frac{\beta}{\alpha(\alpha + \beta M)} \mathbf{1}\mathbf{1}^\top
\end{align}

Công thức này cho phép tính closed-form của $\Theta^{-1}$ mà không cần thuật toán nghịch đảo ma trận $O(M^3)$, đồng thời làm rõ cấu trúc của bộ dự báo NTK.

\section{Các kết quả xác suất}

\subsection{Bổ đề tập trung Gaussian}

Bổ đề tập trung (concentration bound) cho phân phối Gaussian là cơ sở để xác định số lượng mô hình ensemble cần thiết.

\textbf{Bổ đề A.1 (Gaussian Concentration):} Cho $X_1, X_2, \ldots, X_n$ là các biến ngẫu nhiên độc lập, cùng phân phối $\mathcal{N}(\mu, \sigma^2)$. Gọi $\overline{X}_n = \frac{1}{n}\sum_{i=1}^{n} X_i$ là trung bình mẫu. Khi đó, với mọi $t > 0$:
\begin{equation}
    \mathbb{P}\left\{|\overline{X}_n - \mu| \geq t\right\} \leq 2 \exp\left(-\frac{nt^2}{2\sigma^2}\right)
\end{equation}

\textbf{Ứng dụng trong tính ensemble size:}

Đầu ra của mỗi mô hình trong ensemble tuân theo phân phối Gaussian với kỳ vọng $\mu$ (có dấu đúng) và phương sai $\sigma^2$. Để đảm bảo trung bình ensemble $\overline{F}$ có dấu đúng với xác suất ít nhất $1 - \delta$, ta cần:
\begin{equation}
    \mathbb{P}\{|\overline{F} - \mu| \geq |\mu|\} \leq \delta
\end{equation}

Áp dụng Bổ đề A.1 với $t = |\mu|$:
\begin{align}
    2\exp\left(-\frac{N\mu^2}{2\sigma^2}\right) &\leq \delta \\
    N &\geq \frac{2\sigma^2}{\mu^2} \ln\left(\frac{2}{\delta}\right)
\end{align}

Đây chính là Công thức (8) trong main paper, cung cấp ngưỡng chặn trên cho số mô hình cần thiết.

\section{Chi tiết tính toán bộ dự báo NTK}

Phần này trình bày công thức tính các trọng số quyết định trong đầu ra của bộ dự báo NTK.

\subsection{Phân rã đầu ra NTK}

Đầu ra của bộ dự báo NTK tại điểm test $\hat{\mathbf{x}}$ có thể được phân rã:
\begin{equation}
    \mu(\hat{\mathbf{x}})_i = \sum_{j=1}^{M} w_j(\hat{\mathbf{x}}) \cdot y_j
\end{equation}

trong đó $w_j(\hat{\mathbf{x}}) = [\Theta(\hat{\mathbf{x}}, X) \cdot \Theta(X, X)^{-1}]_j$ là trọng số đóng góp của mẫu huấn luyện thứ $j$.

\subsection{Trọng số tín hiệu và nhiễu}

Với cấu trúc block-partitioned của dữ liệu, các trọng số được phân loại thành hai nhóm:

\textbf{Trọng số tín hiệu} $w^1(\hat{\mathbf{x}})$: Áp dụng cho template khớp chính xác (tích vô hướng bằng 1):
\begin{equation}
    w^1(\hat{\mathbf{x}}) = \frac{1}{\alpha} - \frac{\beta}{(\alpha + \beta M)\alpha} \cdot \sum_{j \in \text{match}} y_j
\end{equation}

\textbf{Trọng số nhiễu} $w^0(\hat{\mathbf{x}})$: Áp dụng cho các template có tương quan không hoàn toàn (tích vô hướng khác 0 và 1):
\begin{equation}
    w^0(\hat{\mathbf{x}}) = \frac{\gamma}{\alpha} - \frac{\beta \gamma}{(\alpha + \beta M)\alpha}
\end{equation}

trong đó $\gamma$ là giá trị tích vô hướng giữa $\hat{\mathbf{x}}$ và template tương ứng.

\subsection{Điều kiện để dấu đúng}

Đầu ra có dấu đúng khi thành phần tín hiệu lấn át thành phần nhiễu:
\begin{equation}
    |w^1 \cdot y^*| > \left|\sum_{j \in U} w^0_j \cdot y_j\right|
\end{equation}

trong đó $U$ là tập các template có tương quan không mong muốn.

\section{Phân tích tiệm cận}

\subsection{Lemma 6.1: Bậc của phương sai và kỳ vọng}

Phân tích tiệm cận (\textit{asymptotic orders}) cho phép xác định cách ensemble complexity tăng theo kích thước bài toán.

\textbf{Lemma 6.1:} Trong chế độ NTK với dữ liệu block-partitioned, với $k'$ là kích thước đầu vào (số thành phần khác không):
\begin{align}
    \sigma^2(\hat{\mathbf{x}}) &\in \mathcal{O}\left(\frac{1}{k'}\right) \\
    \mu(\hat{\mathbf{x}}) &\in \mathcal{O}\left(\frac{1}{k'}\right)
\end{align}

\textbf{Hệ quả về tỷ lệ:}
\begin{equation}
    \frac{\sigma^2(\hat{\mathbf{x}})}{\mu^2(\hat{\mathbf{x}})} \in \mathcal{O}(k')
\end{equation}

\textbf{Phân tích theo số block khớp $n_{\hat{x}}$:}

Khi đầu vào $\hat{\mathbf{x}}$ khớp với $n_{\hat{x}}$ template (mỗi template ở một block khác nhau):
\begin{itemize}
    \item \textbf{Kỳ vọng:} $\mu \propto n_{\hat{x}} \cdot w^1$, tăng tuyến tính theo số block khớp
    \item \textbf{Phương sai:} $\sigma^2 \propto n_{\hat{x}} \cdot (w^1)^2$, cũng tăng tuyến tính
    \item \textbf{Tỷ lệ:} $\sigma^2/\mu^2 \propto 1/n_{\hat{x}}$, giảm khi số block tăng
\end{itemize}

Kết quả này giải thích tại sao các bài toán với đầu vào sparse hơn (ít block khớp hơn) yêu cầu ensemble lớn hơn.

\chapter{Mã nguồn cài đặt mô phỏng}

Phụ lục này mô tả cấu trúc mã nguồn và quy trình thực thi mô phỏng để xác minh các kết quả lý thuyết.

\section{Môi trường và thư viện}

\subsection{Thư viện Neural Tangents}

Nghiên cứu sử dụng thư viện \textit{Neural Tangents} (Google Research) để tính toán NTK và dự báo đầu ra của mạng nơ-ron trong giới hạn chiều rộng vô hạn.

\textbf{Đặc điểm chính:}
\begin{itemize}
    \item Tính toán closed-form của kernel NTK cho các kiến trúc mạng phổ biến
    \item Hỗ trợ \textit{NTK predictor} để dự báo đầu ra mà không cần huấn luyện thực tế
    \item Tích hợp với hệ sinh thái JAX để tối ưu hóa hiệu năng
\end{itemize}

\subsection{Nền tảng JAX}

\textit{JAX} (Google) được sử dụng làm backend cho các tính toán ma trận hiệu năng cao:
\begin{itemize}
    \item Hỗ trợ tự động vi phân (automatic differentiation) cho tính toán gradient
    \item Tối ưu hóa JIT compilation cho tốc độ thực thi
    \item Hỗ trợ GPU/TPU để xử lý ma trận lớn
\end{itemize}

\textbf{Cài đặt môi trường:}
\begin{lstlisting}
pip install jax jaxlib neural-tangents
\end{lstlisting}

\section{Cấu trúc mã nguồn}

Toàn bộ mã nguồn được công khai tại GitHub:

\begin{center}
\url{https://github.com/BuhDuy256/add_multiply_execute_algo_instructions_with_neural_networks}
\end{center}

Repository được tổ chức thành các module chính:

\subsection{Modules định nghĩa Template}

Mỗi thuật toán có module riêng định nghĩa các mẫu (templates):

\begin{itemize}
    \item \texttt{templates/addition.py}: Định nghĩa 8 template cho mỗi vị trí bit của bộ cộng ripple-carry
    \item \texttt{templates/multiplication.py}: Template cho thuật toán shift-and-add
    \item \texttt{templates/sbn.py}: Template cho chỉ thị Subtract and Branch if Negative
\end{itemize}

\subsection{Scripts mã hóa dữ liệu}

Module \texttt{encoding/} chịu trách nhiệm mã hóa trực chuẩn (\textit{orthonormal encoding}):
\begin{itemize}
    \item Chuyển đổi template thành vector block-partitioned
    \item Đảm bảo tính trực giao giữa các block
    \item Chuẩn hóa vector để có norm phù hợp
\end{itemize}

\subsection{Scripts tính toán Ensemble}

Module \texttt{ensemble/} thực hiện:
\begin{itemize}
    \item Tính toán ensemble mean từ nhiều mô hình
    \item Áp dụng hàm rounding để chuyển đầu ra về $\{-1, +1\}$
    \item Đánh giá độ chính xác trên tập kiểm tra
\end{itemize}

\section{Quy trình thực thi mô phỏng}

Quy trình mô phỏng gồm bốn bước chính:

\subsection{Bước 1: Khởi tạo cấu trúc Block}

Với thuật toán hoạt động trên $l$ bit, khởi tạo cấu trúc block:
\begin{lstlisting}
def initialize_blocks(l, algorithm_type):
    """
    Khoi tao cau truc block cho thuat toan
    Args:
        l: So bit dau vao
        algorithm_type: 'addition', 'multiplication', 'sbn'
    Returns:
        block_config: Cau hinh block cho tung vi tri
    """
    if algorithm_type == 'addition':
        n_blocks = l  # Moi vi tri bit mot block
        block_sizes = [3] * l  # 3 bit dau vao moi block
    # ...
    return block_config
\end{lstlisting}

\subsection{Bước 2: Tạo tập dữ liệu huấn luyện}

Sinh tất cả các template từ cấu hình block:
\begin{lstlisting}
def generate_training_set(block_config):
    """
    Tao tap du lieu huan luyen tu cac template
    Returns:
        X: Ma tran dau vao (M x d)
        Y: Ma tran nhan (M x output_dim)
    """
    templates = []
    for block_id, config in enumerate(block_config):
        block_templates = enumerate_block_templates(config)
        encoded = encode_block_partitioned(block_templates, block_id)
        templates.extend(encoded)
    return np.array(templates)
\end{lstlisting}

\subsection{Bước 3: Tính toán ma trận NTK}

Sử dụng Neural Tangents để tính kernel:
\begin{lstlisting}
from neural_tangents import stax

# Dinh nghia kien truc mang
init_fn, apply_fn, kernel_fn = stax.serial(
    stax.Dense(n_hidden),
    stax.Relu(),
    stax.Dense(output_dim)
)

# Tinh NTK tren tap huan luyen
K_train = kernel_fn(X_train, X_train, 'ntk')
K_test_train = kernel_fn(X_test, X_train, 'ntk')

# Tinh NTK predictor
predictions = K_test_train @ np.linalg.solve(K_train, Y_train)
\end{lstlisting}

\subsection{Bước 4: Thực thi lặp và làm tròn}

Với các thuật toán đa bước (multi-step), thực thi lặp:
\begin{lstlisting}
def iterative_execution(x_init, n_steps, predictor):
    """
    Thuc thi lap thuat toan
    Args:
        x_init: Trang thai ban dau
        n_steps: So buoc lap
        predictor: Bo du bao NTK
    Returns:
        x_final: Trang thai ket thuc
    """
    x = x_init
    for step in range(n_steps):
        # Du bao buoc tiep theo
        x_pred = predictor(x)
        # Lam tron ve {-1, +1}
        x = np.sign(x_pred)
    return x
\end{lstlisting}

\section{Kiểm chứng thực nghiệm}

\subsection{Trường hợp kiểm tra}

Nhóm nghiên cứu đã xác minh tính đúng đắn trên các trường hợp cụ thể:

\textbf{Phép cộng 10-bit:}
\begin{itemize}
    \item Kiểm tra tất cả $2^{20}$ cặp đầu vào (hai số 10-bit)
    \item Độ chính xác đạt 100\% sau khi làm tròn
    \item Thời gian tính toán: dưới 5 phút trên GPU
\end{itemize}

\textbf{Phép nhân 8-bit:}
\begin{itemize}
    \item Kiểm tra ngẫu nhiên 10,000 cặp đầu vào
    \item Độ chính xác đạt 100\% cho tất cả các trường hợp
    \item Xác nhận tính đúng đắn của cơ chế thực thi lặp
\end{itemize}

\subsection{Kết quả hội tụ}

Các thực nghiệm ghi nhận sự hội tụ tuyệt đối (exact convergence):
\begin{itemize}
    \item Mọi bit đầu ra đều có dấu đúng sau làm tròn
    \item Không có accumulated error qua các bước lặp
    \item Kết quả khớp hoàn toàn với tính toán số học tiêu chuẩn
\end{itemize}

Kết quả này xác nhận tính đúng đắn của cả khung lý thuyết lẫn triển khai mã nguồn.
